\documentclass[12pt]{article}

\usepackage[margin=1in]{geometry}
\usepackage{amsmath,amssymb,amsthm}
\usepackage{mathtools}
\usepackage{booktabs}
\usepackage{algorithm}
\usepackage{algpseudocode}
\usepackage{hyperref}
\usepackage[capitalize]{cleveref}
\usepackage{lineno}
\linenumbers

\newtheorem{proposition}{Proposition}

\newcommand{\Tsys}{T_{\mathrm{sys}}}
\newcommand{\fsys}{f_{\mathrm{sys}}}
\newcommand{\btheta}{\boldsymbol{\theta}}
\newcommand{\blambda}{\boldsymbol{\lambda}}
\newcommand{\bx}{\mathbf{x}}

\title{Supplementary Material for\\
  ``A Binary Threshold for Component Identification\\
  in $k$-out-of-$m$ Systems''}
\author{Alexander Towell}
\date{}

\begin{document}
\maketitle

\appendix

\section{Proof of Proposition 1 (Permutation Invariance)}
\label{sec:proof-symmetry}

We provide the full proof of the permutation invariance of the system
density for $k$-out-of-$m$ systems.

\begin{proof}
Let $\pi$ be a permutation of $\{1, \ldots, m\}$.  Define
$\blambda_\pi = (\lambda_{\pi(1)}, \ldots, \lambda_{\pi(m)})$.  Under
the permuted parameters, component $j$ has rate $\lambda_{\pi(j)}$.
The system density under $\blambda_\pi$ is
\begin{equation}\label{eq:proof-step1}
  \fsys(t; \blambda_\pi) = \sum_{j=1}^{m} f_j(t; \lambda_{\pi(j)})
    \sum_{\substack{\bx_{-j}: j \text{ critical}}}
    \prod_{\ell \neq j} p_\ell(x_\ell, t; \lambda_{\pi(\ell)}).
\end{equation}

For a $k$-out-of-$m$ system, the critical states for component $j$ depend
only on the \emph{number} of other components that have failed, not on
their identities: component $j$ is critical precisely when $k - 1$ of the
remaining $m - 1$ components have failed.  The set of critical state
vectors $\bx_{-j}$ is therefore the set of all binary vectors of length
$m - 1$ with exactly $k - 1$ ones.

Now substitute $\ell' = \pi^{-1}(\ell)$ in the outer sum and rearrange
the inner products.  Since the collection of all $(k-1)$-subsets of
$\{1, \ldots, m\} \setminus \{j\}$ is mapped bijectively to the
collection of all $(k-1)$-subsets of $\{1, \ldots, m\} \setminus
\{\pi^{-1}(j)\}$ under $\pi^{-1}$, and since each product term depends
only on the rate assigned to that component (not the component label),
the sum \eqref{eq:proof-step1} is a rearrangement of the density
evaluated at $\blambda$.

Concretely, letting $j' = \pi^{-1}(j)$:
\[
  \fsys(t; \blambda_\pi)
  = \sum_{j'=1}^{m} f_{j'}(t; \lambda_{j'})
    \sum_{\substack{\bx_{-j'}: j' \text{ critical}}}
    \prod_{\ell' \neq j'} p_{\ell'}(x_{\ell'}, t; \lambda_{\ell'})
  = \fsys(t; \blambda).
\]
The first equality uses the fact that $j = \pi(j')$ is a bijection on
$\{1, \ldots, m\}$ and that the critical states for $\pi(j')$ under
$\blambda_\pi$ correspond to the critical states for $j'$ under
$\blambda$.  The second equality is the definition of $\fsys$.
\end{proof}


\section{Inclusion-Exclusion Expansion}
\label{sec:ie-expansion}

For exponential parallel systems ($k = m$), the product $\prod_{i \neq j}
(1 - e^{-\lambda_i t})$ can be expanded via inclusion-exclusion:
\begin{equation}\label{eq:ie}
  \prod_{i \neq j} (1 - e^{-\lambda_i t})
  = \sum_{S \subseteq \{1,\ldots,m\} \setminus \{j\}}
    (-1)^{|S|} \exp\Bigl(-\sum_{i \in S} \lambda_i \cdot t\Bigr).
\end{equation}
This has $2^{m-1}$ terms, each a pure exponential in $t$.  The component
contribution $w_j(t) = \lambda_j e^{-\lambda_j t} \prod_{i \neq j}
(1 - e^{-\lambda_i t})$ is therefore a signed sum of $2^{m-1}$
exponentials:
\[
  w_j(t) = \lambda_j \sum_{S \subseteq \{1,\ldots,m\} \setminus \{j\}}
    (-1)^{|S|} \exp\bigl(-(\lambda_j + \textstyle\sum_{i \in S} \lambda_i)
    \cdot t\bigr).
\]

This representation makes all integrals closed-form:
\[
  \int_a^b w_j(t)\, dt = \lambda_j
    \sum_S (-1)^{|S|} \frac{e^{-r_S a} - e^{-r_S b}}{r_S},
\]
where $r_S = \lambda_j + \sum_{i \in S} \lambda_i$.  This is used for
interval-censored and left-censored observations in the Scheme~0 and
Scheme~1 likelihoods.  The expansion has $O(2^m)$ terms and is practical
for $m \leq 15$.


\section{Partial Autopsy Enumeration}
\label{sec:autopsy-algorithm}

\begin{algorithm}[t]
\caption{Partial autopsy likelihood evaluation for one observation.}
\label{alg:autopsy}
\begin{algorithmic}[1]
\Require{System time $t$, inspected set $\mathcal{I}$, observed failed
  set $F_{\text{obs}}$, observed survived set $S_{\text{obs}}$,
  parameters $\btheta$, system parameter $k$}
\State $\mathcal{U} \gets \{1,\ldots,m\} \setminus \mathcal{I}$
  \Comment{uninspected components}
\State $k_{\text{unk}} \gets k - |F_{\text{obs}}|$
  \Comment{failures needed from uninspected}
\If{$k_{\text{unk}} < 0$ \textbf{or} $k_{\text{unk}} > |\mathcal{U}|$}
  \State \Return $-\infty$ \Comment{inconsistent observation}
\EndIf
\State $L \gets 0$
\For{each $U \in \binom{\mathcal{U}}{k_\text{unk}}$}
  \Comment{enumerate unknown failed subsets}
  \State $F \gets F_{\text{obs}} \cup U$
  \Comment{full failed set}
  \For{each $j \in F$}
    \Comment{critical component}
    \State $\text{term} \gets f_j(t) \cdot
      \prod_{\ell \in F \setminus \{j\}} F_\ell(t) \cdot
      \prod_{\ell \notin F} S_\ell(t)$
    \State $L \gets L + \text{term}$
  \EndFor
\EndFor
\State \Return $\log L$
\end{algorithmic}
\end{algorithm}

Algorithm~\ref{alg:autopsy} gives the pseudocode for evaluating the
partial autopsy likelihood for a single observation.  The outer loop
enumerates all $\binom{m-r}{k_{\text{unk}}}$ possible assignments of the
unknown failures to uninspected components.  The inner loop over $j \in F$
computes the system density contribution for each candidate critical
component.  Total cost per observation is
$O(\binom{m-r}{k_{\text{unk}}} \cdot k)$.

\end{document}
